\section{Derivations}
\label{app:proofs}

\paragraph{Proposition~\ref{prop:memo}.}
With non-repeating contexts, the empirical loss of a conditional generator decomposes over training contexts. For the memorizing solution every candidate equals the realized mode, so the minimum distance is zero for every context. For the generalizing solution the candidates are $m$ i.i.d.\ draws from $p(\cdot\,|\,c_i)$, and the realized mode is itself a draw from $p$, so the expected minimum distance is $I(m)$ of Equation~\ref{eq:I}. The difference of the two empirical risks is therefore $I(m)$. For the bounds, let the modes be equiprobable and separated by distances in $[\delta, \Delta]$. The probability that none of the $m$ draws equals the realized mode $k$ is $(1 - 1/K)^m$; on that event the nearest draw is at squared distance in $[\delta^2, \Delta^2]$, and otherwise the distance is zero, which gives the stated bounds. If the modes densely fill a curve of length $\ell$ ($K \gg m$), the draws are $m$ uniform points on the curve and the expected squared arc distance from a uniform point to its nearest draw is $\Theta(\ell^2/m^2)$, as for the nearest of $m$ uniform points on an interval, so $I(m)$ decays polynomially and does not depend on $K$. For squared-error regression the single prediction can memorize the realized mode at zero loss, and the generalizing prediction (the conditional mean) has loss $\mathrm{tr}\,\mathrm{Cov}_p$, independent of any budget. A detail used in Section~\ref{sec:collapse}: the realized-mode share $s$ relates to the excess memorization through $s = 1/K + (1 - 1/K)\,\mathrm{excess}$, so a $\pm 2$ RMSE interval on the excess scale maps to $\pm 2(1 - 1/K)\,\mathrm{RMSE}$ on the index scale.

\paragraph{Proposition~\ref{prop:shrink}.}
Write $Y_j - X = (\mu - X) + s\,\varepsilon_j$ with $\varepsilon_j \sim \mathcal{N}(0, I_d)$ and $X - \mu \sim \mathcal{N}(0, \sigma^2 I_d)$. Then $\|Y_j - X\|^2 = \|X - \mu\|^2 + s^2 \|\varepsilon_j\|^2 - 2 s\, \varepsilon_j^\top (X - \mu)$. The first term is common to all $j$; the second concentrates at $s^2 d$ with fluctuations of order $s^2 \sqrt{d}$; the third is, conditionally on $X$, Gaussian with variance $4 s^2 \|X - \mu\|^2 \approx 4 s^2 \sigma^2 d$ and independent across $j$, so the minimum over $j$ of the sum is, to leading order in $d$, $d\sigma^2 + d s^2 - 2 s \sigma \sqrt{d}\, e_m$ with $e_m = \E \max_{j \le m} \xi_j$. Minimizing in $s$ gives $s^\star = \sigma e_m / \sqrt{d}$. The spread-to-error ratio used in Section~\ref{sec:bair}, defined as the mean pairwise distance among samples divided by the mean distance from samples to the truth, is $r = \sqrt{2}\rho/\sqrt{1 + \rho^2}$ with $\rho = s/\sigma$ under this model, which is how $\deff$ is fitted from one measured $r$.

\paragraph{Flattening.}
Let the generator place mass $q_k$ on mode $k$ and let all modes be at mutual distance $\Delta$. The probability that none of $m$ draws lands on the realized mode $k$ is $(1 - q_k)^m$, so $L(q) = \Delta^2 \sum_k p_k (1 - q_k)^m$. Minimizing subject to $\sum_k q_k = 1$ gives the stationarity condition $m\, p_k (1 - q_k)^{m-1} = \lambda$ for every $k$ with $q_k > 0$, which is Theorem~7 of \citet{bennett2026freedom}; for $K = 2$ it reads $\big((1-q)/q\big)^{m-1} = (1-p)/p$. If $p_k = 1/K$ for all $k$ the symmetric point $q_k = 1/K$ satisfies the condition and is the unique minimizer of the convex objective for $m \ge 2$ (Theorem~4 of the same work). The numerical optima quoted in the text (for $p = (0.787, 0.213)$: $0.787$, $0.607$, $0.522$ at $m = 2, 4, 16$) solve this equation.

\paragraph{Proposition~\ref{prop:es}.}
(a) The energy distance between $q$ and $\delta_x$ is $2\E_{Y \sim q}\|Y - x\|^\beta - \E_{Y, Y' \sim q}\|Y - Y'\|^\beta - 0$, so $\IS_\beta(q, x)$ is half of it; the energy distance with $\beta \in (0, 2)$ is a metric on distributions with finite moments \citep{gneiting2007strictly,dawid2007geometry}, hence non-negative and zero only when $q = \delta_x$. The empirical risk over training contexts is a sum of such terms with one $x$ per context, so its unique minimizer is the memorizing solution; the $m$-sample estimator with the unbiased pairwise term has this risk as its expectation for every $m \ge 2$. (b) At $q = p$, $\E_{x \sim p}\,\IS_\beta(p, x) = \E\|Y - x\|^\beta - \frac{1}{2}\E\|Y - Y'\|^\beta = \frac{1}{2}\E\|Y - Y'\|^\beta$ because $x$ and $Y$ are exchangeable. (c) For $\beta = 2$, $\E\|Y - x\|^2 - \frac{1}{2}\E\|Y - Y'\|^2 = \|\E Y - x\|^2$. (d) The per-token score is a sum over tokens of energy scores of the token marginals, and (a)--(b) apply to each term.

\paragraph{Direction insensitivity (interpretation).}
Suppose the generator's mean misses the outcome by an unknown vector of length $a$ and it adds isotropic noise of scale $s$ in $d$ dimensions with $d$ large. Then $\E\|Y - x\|^\beta \approx (a^2 + s^2)^{\beta/2}$ and $\E\|Y - Y'\|^\beta \approx (2 s^2)^{\beta/2}$, and the energy score $(a^2 + s^2)^{\beta/2} - \frac{1}{2}(2s^2)^{\beta/2}$ is minimized at $s = a$ for every $\beta \in (0, 2)$, where it equals $2^{\beta/2 - 1} a^\beta$; this is the $\lambda = 1$ case of Proposition~4.1 of \citet{debortoli2025distributional}, and the gain does not depend on $d$. The same calculation for the \nom{} loss follows Proposition~\ref{prop:shrink} and gives $s \approx a\, e_m/\sqrt{d}$. The approximation requires the bias and the noise to be spread over many dimensions (participation ratio $\gg 1$); the regression residual on BAIR has participation ratio about $1{,}030$, and we did not measure it on the bounce tokens. The argument says nothing about where the isotropic spread points; the claim in Section~\ref{sec:es} that it points off the manifold is an empirical observation.

\section{$I(m,K)$ for the bounce geometry and the recorded predictions}
\label{app:pred}
\label{app:imk}
At a wall frame the $K$ possible next directions are the unit vectors at angles $\theta_k = \big((k + \tfrac12)/K\big)\,180^\circ - 90^\circ$ about the wall normal, $k = 0, \dots, K-1$, with unit speed. For a target mode $k$ let $d^{(k)}_{(1)} \le \dots \le d^{(k)}_{(K)}$ be the sorted squared distances from $\mu_k$ to all modes ($d^{(k)}_{(1)} = 0$). The nearest of $m$ i.i.d.\ uniform draws is the $r$-th closest mode with probability $\big((K - r + 1)/K\big)^m - \big((K - r)/K\big)^m$, so
\begin{equation}
I(m, K) = \frac{1}{K} \sum_{k=1}^{K} \sum_{r=1}^{K} d^{(k)}_{(r)} \Big[ \big(\tfrac{K - r + 1}{K}\big)^m - \big(\tfrac{K - r}{K}\big)^m \Big].
\end{equation}
Table~\ref{tab:predfull} lists the values together with the predictions recorded before the $K \in \{2, 16\}$ runs. The predictions were made by fitting excess memorization linearly in $I$ on the eight $K \in \{4, 8\}$ points (excess $= 0.015 + 1.06\,I$, RMSE $0.032$) and mapping to the index scale; the intervals are $\pm 2$ RMSE. We noted two risks at the time: the $K{=}2$, $m{=}2$ point lies outside the fitted range of $I$, and large $K$ might memorize less than predicted because more bits must be stored per context. Both misses are of these kinds.

\begin{table}[h]
\caption{Recorded predictions for the memorization index at 8k steps and the observed two-seed means. The two misses are marked.}
\label{tab:predfull}
\begin{center}
\small
\resizebox{\linewidth}{!}{%
\begin{tabular}{cccccc}
\toprule
$K$ & $m$ & $I(m,K)$ & Chance & Predicted [interval] & Observed \\
\midrule
2 & 2 & 0.500 & 0.50 & 0.77 [0.74, 0.81] & 0.956 (miss) \\
2 & 4 & 0.125 & 0.50 & 0.57 [0.54, 0.61] & 0.568 \\
2 & 8 & 0.008 & 0.50 & 0.51 [0.50, 0.54] & 0.528 \\
2 & 16 & $\approx 0$ & 0.50 & 0.51 [0.50, 0.54] & 0.522 \\
16 & 4 & 0.234 & 0.0625 & 0.31 [0.25, 0.37] & 0.221 (miss) \\
16 & 8 & 0.074 & 0.0625 & 0.15 [0.09, 0.21] & 0.144 \\
16 & 16 & 0.021 & 0.0625 & 0.10 [0.06, 0.16] & 0.115 \\
16 & 32 & 0.005 & 0.0625 & 0.08 [0.06, 0.14] & 0.105 \\
16 & 64 & 0.0006 & 0.0625 & 0.08 [0.06, 0.14] & 0.081 \\
\midrule
4 & 4 / 8 / 16 / 32 & 0.235 / 0.061 / 0.006 / 0.0001 & 0.25 & (fitted) & 0.498 / 0.306 / 0.272 / 0.263 \\
8 & 4 / 8 / 16 / 32 & 0.236 / 0.076 / 0.019 / 0.002 & 0.125 & (fitted) & 0.305 / 0.198 / 0.154 / 0.144 \\
\bottomrule
\end{tabular}}
\end{center}
\end{table}

\section{Benchmark construction}
\label{app:bench}

\paragraph{A2 appearance randomization.}
Each episode of A2-div draws a 22-dimensional appearance vector: background RGB in $[10, 110]$, ball RGB in $[150, 255]$, radius in $[3, 6]$\,px, speed in $[10, 18]$\,px per frame, and two rectangles with corner in $[0, 48]$, size in $[6, 20]$ and RGB in $[40, 200]$. Frames are rendered at $64\times64$ with an anti-aliased disc, upsampled to $224$\,px and encoded by DINOv2 ViT-S/14. Context novelty is measured as the ratio of the median distance from a test context to its nearest training context to the median nearest-neighbour distance within the training set in a PCA-64 token space: $1.9$ for A2-div and $0.35$ for A2-fix. The evaluation uses 400 held-out wall contexts; a stochastic transition is one whose current frame touches a wall.

\paragraph{Push environment.}
The $64\times64$ arena contains an ice rectangle ($x \in [20, 44]$, $y \in [14, 40]$) and an absorbing pit ($x \in [24, 40]$, $y \in [0, 14]$). A push action is a 2-D vector; with probability $0.1$ it is zero, otherwise its direction is uniform and its magnitude in $[0.3, 1]$. The ball moves $10$\,px along the action rotated by the sampled deflection and respawns after four frames in the pit. The branch-risk task draws starts with $x \in [8, 12]$, $y \in [17, 21]$ and goals with $x \in [52, 56]$ in the same band. The planner uses CEM-MPC with horizon $T{=}5$, 64 candidate sequences, $m_p{=}8$ imagined futures per candidate (one for regression), and a cost equal to the mean distance between the decoded final position and the goal; 100 episodes per model and seed with a shared set of task seeds. The risk diagnostic evaluates 200 start states with the straight five-push action sequence against 256 true-simulator rollouts per start, and a one-push test from 200 ice states just below the pit.

\paragraph{Stochastic Push-T hidden variable.}
Table~\ref{tab:cog} shows the per-episode randomizations we tried before choosing the centre of gravity; mass and friction have no effect under Push-T's quasi-static physics, and damping produces unrealistic sliding. The multi-future test set takes 128 starting states just before the first contact (an earlier version with earlier starts did not diverge within 8 steps and was discarded), and for each the 64 futures share the three context frames. Modes closer than $2$\,px in outcome are merged, leaving $3.9$ distinguishable modes per start on average. The CEM planner for Push-T uses $T{=}5$, 64 candidates, $m_p{=}4$, 3 iterations with 8 elites, re-planning every step for at most 10 steps, with the mean token-space $\ell_2$ distance to the goal as the cost, following \citet{terver2025drives}; goals are within 25 steps of the start, and success requires the block within $20$\,px and $20^\circ$ of the goal.

\begin{table}[h]
\caption{Candidate hidden variables for stochastic Push-T: deviation of the T vertices after 16 steps under the same actions (median / maximum over episodes; the nominal displacement is $124$\,px).}
\label{tab:cog}
\begin{center}
\small
\begin{tabular}{lcl}
\toprule
Randomization & Deviation (px) & Decision \\
\midrule
Block mass $0.5$--$2\times$ & 0.0 / 0.0 & no effect (quasi-static) \\
Contact friction $0.3$--$1.2$ & 0.0 / 0.0 & no effect \\
Damping $0$--$0.1$ & 25 / 346 & unrealistic sliding \\
Centre of gravity, radius $20$\,px, 4 positions + 3\,px jitter & 28 / 138 & adopted \\
Centre of gravity, radius $25$\,px & 56 / 195 & too large \\
\bottomrule
\end{tabular}
\end{center}
\end{table}

\paragraph{BAIR.}
We use the first part of the official archive: $4{,}096$ training sequences and $1{,}024$ held-out sequences (the official test set plus three held-out training files), $64\times64$ frames, 4-dimensional actions, 3 context frames and 16-step rollouts.

\section{Implementation details and ablations}
\label{app:impl}

\paragraph{Training.}
AdamW with learning rate $2\times10^{-4}$, $(\beta_1, \beta_2) = (0.9, 0.9)$, no weight decay, a linear warm-up over 1{,}000 steps followed by a cosine decay, batch 128, gradient clipping at norm 1, an exponential moving average of the weights with decay $\min(0.999, (1 + t)/(10 + t))$ used for all evaluations, 12k steps (8k for the budget grid, the asymmetric variant and the per-step-matching comparisons; 36k for iMF; 24k for Shortcut; 48k for the four-times-longer A1 runs). The noise vector has 256 dimensions and is normalized to unit RMS before the mapping network. Late fusion sets the first five of eight blocks (eight of twelve for the 21M model) to a learned constant style. Candidate generation for matching runs under no-gradient in bfloat16 in a separate autocast region. The Geman--McClure scale is twice the running median (EMA $0.99$) of the matching distance. The energy-score runs use $\beta = 1$, the whole-map RMS distance, the unbiased pairwise estimator, and backpropagate in two chunks of contexts; the per-token variant computes pairwise differences explicitly to be exact near zero. Flow matching uses the time shift $\alpha = 13$ at training and sampling and the Euler sampler. The iMF port follows the official JAX implementation through a PyTorch re-implementation: dual $u$/$v$ heads of width 640 (to match parameters), the JVP direction taken from the model's own $v$ prediction, and the loss $\|u + (t - r)\,\mathrm{sg}(\partial_t u) - v\|^2$ plus a flow-matching loss on $v$; it converges only with the time-embedding frequency scale set to 1 rather than the usual 1000. Our own MeanFlow implementation diverged three times once the average-velocity term was switched on and was abandoned. Shortcut models devote a quarter of each batch to the self-consistency target computed with the EMA model. Parameter counts: 8M setting, $7.97$--$8.24$M across methods ($8.24$M for \imlewm{}); 21M setting, $20.6$M for \imlewm{}. Training times in the tables are wall-clock hours on one RTX A5000; some runs shared a GPU and are slower than the stand-alone figures (8M, batch 128, one step: $m{=}4$ 178\,ms, $m{=}16$ 488\,ms; 21M: 404 and 1112\,ms).

\paragraph{Recipe choices.}
Table~\ref{tab:recipe} gives the pilot comparisons behind the recipe. The wide head does not change \imlewm{} (it is there to match FM); late fusion halves training time at equal coverage and lowers the deterministic error; the FlowWM recipe raised the flow baseline substantially. Multi-stage supervision, which ROMS-IMLE found important for images, does not help on a fixed $16\times16$ token grid (Appendix~\ref{app:negative}).

\begin{table}[h]
\caption{Recipe comparisons on A2-div ($K{=}4$, pilot and preparation runs, one seed; bootstrap 95\% intervals over 400 contexts where given). Training hours are on a shared GPU.}
\label{tab:recipe}
\begin{center}
\small
\resizebox{\linewidth}{!}{%
\begin{tabular}{lcccccc}
\toprule
Variant & Coverage & Precision & kNN recall & Det.\ err.\ (px) & Memo.\ (train / val) & Train (h) \\
\midrule
\imlewm{} $m{=}16$, full noise injection, wide head & 0.934 [0.917, 0.949] & 0.83 & 0.94 & 2.31 & 0.29 / 0.24 & 1.93 \\
\imlewm{} $m{=}16$, late fusion, wide head (adopted) & 0.925 [0.907, 0.942] & 0.83 & 0.95 & 1.56 & 0.29 / 0.27 & 1.05 \\
\imlewm{} $m{=}16$, full injection, narrow head (6.6M) & 0.930 [0.909, 0.948] & 0.84 & 0.94 & 2.02 & 0.30 / 0.25 & 1.71 \\
\imlewm{} $m{=}16$, pilot recipe (three spatial stages) & 0.91 [0.89, 0.93] & 0.81 & 0.94 & 2.67 & 0.30 / 0.22 & 3.2 \\
FM, default recipe, 20 Euler steps & 0.58 [0.55, 0.61] & 0.38 & 0.68 & 4.64 & -- & 0.6 \\
FM, FlowWM recipe, 50 steps (adopted) & 0.79 [0.77, 0.82] & 0.51 & 0.91 & 2.27 & 0.29 / 0.24 & 0.4 \\
iMF, time-embedding scale 1000 & \multicolumn{5}{c}{does not converge} \\
iMF, time-embedding scale 1, 12k steps & 0.55 [0.52, 0.58] & 0.22 & -- & 13.1 & -- & 1.0 \\
\bottomrule
\end{tabular}}
\end{center}
\end{table}

\paragraph{Ablations.}
\label{app:ablation}
On A2-div with $m{=}16$ (one seed): replacing Geman--McClure by $\ell_2$ gives coverage $0.93$, precision $0.83$, index $0.28$, deterministic error $1.81$\,px (vs.\ $1.55$); a 64-dimensional noise vector gives $0.92$, $0.83$, $0.28$, $1.50$; training $1.5\times$ longer (18k steps) gives $0.92$, $0.86$, $0.31$, $1.45$. On stochastic Push-T an 8M \imlewm{} with $m{=}16$ reaches recall errors $24.5$ / $39.0$\,px at $h = 8$ / $16$ against $22.1$ / $37.8$ for the 21M model (three seeds), so model size matters little there. In the pilot with the earlier recipe, Geman--McClure improved precision ($0.92$ vs.\ $0.86$) and token-space coverage ($0.85$ vs.\ $0.66$) over $\ell_2$. The iMF learning curve on A2-div (one seed, checkpoints of a 72k-step run) saturates after 24k steps: coverage $0.56$, $0.57$, $0.58$, $0.58$ and precision $0.24$, $0.25$, $0.26$, $0.26$ at 24k, 42k, 60k and 72k steps.

\section{Readouts}
\label{app:readout}
Position readouts are small networks trained on real tokens: a soft-argmax probe for the ball (test error $0.13$\,px on A2-div, $0.05$\,px with DINOv3 on A1), and for Push-T a probe for agent position, block position and angle (T-vertex error $2.0$\,px). Because a probe might decode generated (out-of-distribution) tokens differently from real ones, the pilot cross-checked it against a nearest-neighbour judge that decodes a token map by looking up the nearest of $12{,}000$ real frames in a PCA-128 space: on the fixed-appearance environment the two judges agreed within $0.01$ for every \imlewm{} run (e.g.\ coverage $0.960$ vs.\ $0.960$ at $K{=}4$) and the nearest-neighbour judge was $0.03$--$0.09$ lower for flow matching ($0.841$ vs.\ $0.814$ with 20 steps), so the probe is, if anything, generous to the baselines; the ranking was identical under both. The FlowWM heatmap readout, ported from the public code, maps the $384$-channel map through a $1\times1$ convolution to 256 channels, GELU, dropout $0.1$ and a $16\times$ transposed convolution to $256\times256$ logits per shape, trained with pixel cross-entropy and feature noise of standard deviation $0.05$ for 20k steps, and reads the argmax; its error on real features is $0.10$\,px (p95 $0.75$). The noise-robust variant is trained with feature noise $0.5$ (floor $1.03$\,px). Table~\ref{tab:readout} gives all readout combinations for A1.

\begin{table}[h]
\caption{Bouncing Shapes F1 error (px) under three readouts, three seeds (energy score two). The last row is the 48k-step run of Section~\ref{sec:a1} (one seed).}
\label{tab:readout}
\begin{center}
\small
\begin{tabular}{lcccccc}
\toprule
 & \multicolumn{3}{c}{Joint 16$\to$16} & \multicolumn{3}{c}{Autoregressive} \\
\cmidrule(lr){2-4}\cmidrule(lr){5-7}
Method & probe & heatmap & heatmap, noise 0.5 & probe & heatmap & heatmap, noise 0.5 \\
\midrule
Regression & 5.14 & 22.48 & 9.67 & 8.09 & 8.41 & 8.53 \\
FM-50 & 3.86 & 12.59 & 15.00 & 12.28 & 14.11 & 18.38 \\
FM-2 & 3.83 & 14.80 & 16.43 & 9.95 & 12.49 & 16.85 \\
Energy score, $m{=}8$ & -- & -- & -- & 8.94 & 33.77 & 11.65 \\
\imlewm{}, $m{=}4$ & 3.30 & 11.73 & 5.92 & 6.61 & 6.83 & 7.04 \\
\imlewm{}, $m{=}16$ & 3.32 & 9.04 & 5.28 & 7.21 & 7.64 & 7.92 \\
\imlewm{}, $m{=}4$, 48k steps & -- & -- & -- & 4.08 & 4.20 & -- \\
\bottomrule
\end{tabular}
\end{center}
\end{table}

\section{Additional results}
\label{app:more}

\subsection{Rollout metrics}
Table~\ref{tab:e3} completes Figure~\ref{fig:manifold} with the best-of-16 readout error and the spread at the last horizon.

\begin{table}[h]
\caption{Rollouts from real contexts (16 samples, three seeds; iMF on Push-T two): manifold distance, best-of-16 readout error (px) and spread (px) at the last horizon ($h{=}32$ on A2, $16$ on A1 and B).}
\label{tab:e3}
\begin{center}
\small
\setlength{\tabcolsep}{4pt}
\begin{tabular}{lccccccccc}
\toprule
 & \multicolumn{3}{c}{A2 bounce, $h{=}32$} & \multicolumn{3}{c}{A1 Bouncing Shapes, $h{=}16$} & \multicolumn{3}{c}{B Push-T, $h{=}16$} \\
\cmidrule(lr){2-4}\cmidrule(lr){5-7}\cmidrule(lr){8-10}
Method & manifold & best-of-16 & spread & manifold & best-of-16 & spread & manifold & best-of-16 & spread \\
\midrule
Regression & 1.80 & 32.2 & 0 & 2.04 & 34.0 & 0 & 3.21 & 102.2 & 0 \\
FM-50 & 6.16 & 9.8 & 25.7 & 2.95 & 28.0 & 25.1 & 4.60 & \best{67.4} & 36.4 \\
FM-2 & 6.36 & 10.6 & 26.8 & 3.03 & 23.4 & 26.1 & 4.37 & 69.2 & 29.7 \\
iMF & 8.04 & 10.1 & 29.0 & 3.57 & 33.4 & -- & 5.81 & 84.1 & 57.0 \\
\imlewm{}, $m{=}4$ & 1.74 & 10.5 & 23.9 & \best{1.98} & \best{19.0} & -- & \best{3.20} & 88.7 & 9.2 \\
\imlewm{}, $m{=}16$ & \best{1.68} & \best{9.3} & 27.2 & 2.08 & 19.4 & 17.9 & 3.28 & 79.1 & 21.8 \\
\bottomrule
\end{tabular}
\end{center}
\end{table}

\subsection{Stochastic Push-T: bootstrap intervals}
Table~\ref{tab:bci} gives the paired differences behind Section~\ref{sec:b} and Section~\ref{sec:seq}.

\begin{table}[h]
\caption{Paired differences on the stochastic Push-T multi-future test set (px; negative favours the first model) with 95\% bootstrap intervals over the 128 starting states after averaging seeds, and hierarchical intervals that also resample the three seeds (10{,}000 resamples).}
\label{tab:bci}
\begin{center}
\small
\begin{tabular}{llcccc}
\toprule
Comparison & Metric & Difference & Start-state CI & Hierarchical CI \\
\midrule
\imlewm{} 16 vs.\ FM-50 & $h{=}8$ recall & $-1.8$ & $[-2.5, -1.1]$ & $[-5.2, +2.3]$ \\
 & $h{=}8$ precision & $-5.2$ & $[-5.7, -4.8]$ & $[-8.0, -2.5]$ \\
 & $h{=}16$ recall & $-4.2$ & $[-5.3, -3.1]$ & $[-8.8, +1.1]$ \\
 & $h{=}16$ precision & $-6.6$ & $[-7.6, -5.7]$ & $[-8.5, -4.7]$ \\
\imlewm{} 16 vs.\ FM-2 & $h{=}16$ recall & $-3.7$ & $[-4.9, -2.5]$ & $[-8.7, +1.9]$ \\
 & $h{=}16$ precision & $-5.6$ & $[-6.6, -4.6]$ & $[-7.7, -3.2]$ \\
\imlewm{} 16 vs.\ Shortcut & $h{=}16$ recall & $-8.2$ & $[-10.0, -6.3]$ & $[-11.3, -4.0]$ \\
 & $h{=}16$ precision & $-13.2$ & $[-15.1, -11.3]$ & $[-16.3, -10.2]$ \\
Per-step vs.\ one-step matching (8k) & $h{=}8$ recall & $-7.1$ & $[-8.2, -6.0]$ & $[-9.2, -4.8]$ \\
 & $h{=}16$ recall & $-14.3$ & $[-16.3, -12.5]$ & $[-19.1, -9.9]$ \\
 & $h{=}16$ precision & $-17.0$ & $[-18.4, -15.6]$ & $[-20.4, -14.5]$ \\
\bottomrule
\end{tabular}
\end{center}
\end{table}

\subsection{Energy-score follow-ups}
\label{app:es}
Figure~\ref{fig:toy} shows the toy experiment of Section~\ref{sec:es}. The task is the two-branch problem of \citet{makabe2026esp} ($u \sim U[0.05, 1]$, $y = \pm\sqrt{u} + 0.08\varepsilon$, equiprobable branches) with a 16-dimensional nuisance vector $a \sim \mathcal{N}(0, I)$ appended to the context, a three-layer MLP generator of width 256 with a 4-dimensional noise input, Adam at $2\times10^{-4}$, batch 256, 15k steps, and the same energy-score estimator as the main runs. With fresh contexts at every step the energy score matches $0.940$ / $0.946$ of its samples to a branch centre within $0.24$ ($m = 8$ / $16$) and \nom{} $0.916$ / $0.924$. With 1{,}024 fixed training contexts (about 3{,}750 revisits each) the energy score memorizes ($0.999$) and covers both branches on $0.2\%$ of new contexts, while \nom{} with $m = 8$ / $16$ is at $0.52$ / $0.51$ and covers $0.99$ / $1.00$. With 4{,}096, 16{,}384 and 65{,}536 contexts the energy score's realized-branch share is $0.73$, $0.51$ and $0.50$; without the nuisance vector (context $= u$ only) it is $0.50$ even with 1{,}024 contexts. The thresholds are those of this small network and do not transfer; the 8M model on A2-div saw each of its $11{,}559$ stochastic training contexts about 38 times and memorized.

\begin{figure}[h]
\centering
\includegraphics[width=0.9\linewidth]{fig_toy.pdf}
\caption{Toy two-branch problem with a nuisance appearance vector (two seeds, $m{=}8$ for the energy score and $16$ for \nom{}). (a) Realized-branch share on training contexts against the number of fixed training contexts. (b) Two-branch coverage on new contexts; ``fresh'' draws a new context at every step.}
\label{fig:toy}
\end{figure}

The follow-up runs on A2-fix and with the per-token score were preceded by five recorded predictions, all of which held (two seeds each): (P1) the energy score with $m{=}8$ does not memorize on A2-fix (index $\le 0.40$; observed $0.268$); (P2) its coverage on A2-fix exceeds its A2-div value but stays below \imlewm{} with $m{=}16$ in both coverage and precision (observed $0.947$ vs.\ $0.987$ and $0.814$ vs.\ $0.954$); (P3) \imlewm{} with $m{=}4$ does not memorize on A2-fix ($\le 0.35$; observed $0.256$); (P4) the per-token energy score memorizes on A2-div (index $\ge 0.7$, coverage $\le 0.5$; observed $0.970$ and $0.224$); (P5) its one-step manifold distance exceeds $1.1$ (observed $1.12$; \imlewm{} $0.86$). On A2-fix the energy score's kNN recall is $1.00$ against $0.85$ for \imlewm{}, its samples spread more widely than the outcomes; its deterministic error is the lowest ($0.20$\,px). On A2-div its learning curves peak in coverage at 6k steps ($0.50$) and decline as the index rises ($0.42 \to 0.87$); the two $m{=}16$ seeds reach an index of $0.84$ and $0.87$ at 8k steps against $0.51$--$0.72$ for the three $m{=}8$ seeds. Token-space diagnostics at $h{=}1$ on A2-div (three seeds, evaluation contexts): spread-to-error ratio $0.80$ (energy score) vs.\ $0.64$ (\imlewm{} 16), $0.38$ (\imlewm{} 4), $1.15$ (FM-50); best-of-16 token error $0.468$ vs.\ $0.271$; mean-prediction error $0.451$ vs.\ $0.396$; readout spread $1.7$ vs.\ $3.9$\,px; one-step manifold distance $1.19$ vs.\ $0.86$.

\subsection{Per-step matching}
\label{app:seq}
Per-seed Push-T recall errors at $h{=}16$ (one-step $\to$ per-step): $43.7 \to 30.3$, $42.6 \to 32.2$, $44.7 \to 25.5$. On the deterministic Push-T control (one seed) the best-of-16 T-vertex error at $h = 4 / 8 / 16$ moves from $35.8 / 60.4 / 72.3$ to $29.2 / 45.6 / 58.9$\,px and the manifold distance at $h{=}16$ from $3.90$ to $3.64$. On A2-div (two seeds, 8k steps) the manifold distance at $h = 8 / 16 / 32$ moves from $1.31 / 1.47 / 1.83$ to $1.19 / 1.26 / 1.40$ while one-step coverage and precision move from $0.92 / 0.77$ to $0.90 / 0.73$ and the deterministic error from $2.91$ to $3.97$\,px; the best-of-16 error is unchanged. On A1 (two seeds) F1 moves from $8.74$ to $8.51$ (precision $9.26 \to 9.77$, recall $8.28 \to 7.55$), the manifold distance at $h{=}16$ from $2.17$ to $1.92$ and the best-of-16 error from $23.2$ to $17.7$\,px. The fed-back-state ablation was preceded by three recorded predictions: (S1) feeding back the true frame yields less than half of the manifold gain (manifold at $h{=}16$ $\ge 2.05$; observed $2.14$); (S2) feeding back a random candidate yields more than half ($\le 2.05$; observed $1.81$); (S3) the random candidate's best-of-16 error exceeds the matched one's by at least $1$\,px ($\ge 18.7$; observed $22.9$). The spread of the 16 rollouts at $h = 1 / 4 / 8 / 16$ is $0.20 / 1.90 / 6.81 / 23.6$\,px for the true futures, $0.70 / 3.35 / 7.42 / 17.0$ for one-step matching, $1.16 / 5.42 / 12.2 / 26.3$ for the matched variant, $2.41 / 4.19 / 6.38 / 12.4$ for the random variant and $0.61 / 3.07 / 7.08 / 16.6$ for the true-frame variant. Training times: one-step $0.68$\,h, per-step variants $2.7$\,h (8M, 8k steps).

\subsection{Planning details}
On stochastic Push-T the final block position and angle errors are $46.6$\,px / $42.0^\circ$ (regression), $40.6$ / $35.2$ (FM-2), $45.2$ / $20.0$ (FM-50, 10 episodes) and $49.5$ / $44.9$ (\imlewm{}). The 1-step contact error (block displacement after a contact transition, true $15.9$\,px) is $11.7$ / $11.2$\,px for regression at 12k / 36k steps, $12.4$ / $12.6$ for FM-2 / FM-50 and $12.2$ / $12.3$ for \imlewm{} with $m = 4$ / $16$ (seed 0). On the branch-risk task the per-start absolute error of the five-push pit probability is $0.078$ (\imlewm{}), $0.086$ (FM-50), $0.096$ (FM-2) and $0.462$ (regression), and the timeout rates are $13.7\%$, $14.3\%$, $19.3\%$ and $17\%$ (oracle). Wall-clock for 100 episodes: oracle $100$\,s, regression $523$\,s ($m_p{=}1$), \imlewm{} $3{,}633$\,s, FM-2 $9{,}067$\,s.

\section{Negative results}
\label{app:negative}
The pilot study tested four additions proposed in our original design; none survived.

\paragraph{Trajectory-level matching.}
Matching $m{=}4$ candidate rollouts of length $H{=}8$ as whole trajectories (with a curriculum $H = 1 \to 4 \to 8$) on the push environment was worse than one-step matching at every horizon: best-of-16 position error $1.31$ vs.\ $0.35$\,px at $h{=}1$, $3.89$ vs.\ $2.19$ at $h{=}8$, $4.67$ vs.\ $3.03$ at $h{=}16$ (paired bootstrap intervals exclude zero), one-step coverage $0.71$ vs.\ $0.98$, and $3.3\times$ the training time; only the error of the sample mean improved ($9.18$ vs.\ $9.94$ at $h{=}16$), the signature of averaging. With two to three branches per step an 8-step trajectory has hundreds of modes, far beyond $m{=}4$, so the matched candidate is the one close to all of them on average. Per-step matching (Section~\ref{sec:seqmatch}) replaced it.

\paragraph{Multi-stage supervision.}
On the fixed $16\times16$ token grid, supervising intermediate heads with spatially pooled targets ($4\times4$, $8\times8$, $16\times16$) or with PCA-truncated targets (32, 128, 384 components) gave the same coverage as a single head within the bootstrap interval ($0.960$ and $0.959$ vs.\ $0.967$ at $K{=}4$) and the same precision ($0.916$, $0.922$ vs.\ $0.926$). The final recipe uses a single head.

\paragraph{Variance-penalized planning.}
A CEM cost of mean plus $\lambda$ times the spread of the imagined goal distances lowered the success rate of the branch-risk task even when the true simulator served as the world model: $78\%$, $77\%$, $70\%$ and $50\%$ at $\lambda = 0, 0.5, 1, 2$, with the pit rate rising to $35\%$. A certain failure has zero variance, so the penalty prefers it to an uncertain success. With \imlewm{} as the model (two seeds, 60 episodes each) the pooled success was $62.5\%$ at $\lambda = 0$ and $61.7\%$ at $\lambda = 0.5$.

\paragraph{Counterfactual action contrast.}
We had anticipated that the noise $z$ might absorb the effect of the action. It does not: on the push environment the counterfactual precision (samples under a random other action staying away from the true outcome) and the action sensitivity (output change per action change relative to the truth) are $0.95$ and $0.96$ with the action on its own FiLM path, $0.92$ and $0.97$ with the action merged into the noise path, and $0.93$ and $0.975$ with an added margin term, against $0.83$ and $0.94$ for regression. There was no shortcut to repair.

\paragraph{Low-dimensional matching and larger budgets on BAIR.}
See Table~\ref{tab:bair}: $m{=}64$ raises the spread-to-error ratio from $0.44$--$0.51$ to $0.53$--$0.55$ at $2.4\times$ the training time, and matching in a 512-dimensional pooled space lowers it to $0.31$--$0.38$; neither approaches FM-50 on distribution metrics.

\section{Statistical methodology}
\label{app:stats}
Per-context intervals are percentile bootstrap intervals (2{,}000 resamples) over the 400 evaluation contexts of a run, used in the pilot and recipe tables. Seed statistics are means and standard deviations over the listed seeds. For the Push-T multi-future test set, each model's per-start errors were stored; the start-state interval averages the seeds first and resamples the 128 starts, and the hierarchical interval resamples the seeds with replacement and then the starts, 10{,}000 times each. Per-step and one-step matching runs with the same seed are paired. All comparisons report the direction and the magnitude of the difference; where an interval includes zero we say so in the text.
